% !PRINTED 21
\weberpage{21}
%<<CONT
sein. Diese Sätze gelten, wenn nur eine der beiden Gruppen $Q$, $Q'$
grösster Normaltheiler von $P$ ist. Wir nehmen jetzt an, dass sie
beide diese Eigenschaft haben, und verstehen unter $R$ den grössten
gemeinschaftlichen Theiler von $Q$ und $Q'$. Diese Gruppe $R$ ist
dann ein Normaltheiler von $P$ sowohl als von $Q$ und $Q'$. Denn ist
$a$ irgend ein Element in $P$, und $c$ in $Q$ sowohl als in $Q'$
enthalten, so ist auch $a^{-1}ca$ in $Q$ und in $Q'$, also auch in
$R$ enthalten. Also ist $a^{-1}Ra = R$, und $R$ ist Normaltheiler
von $P$ und mithin Normaltheiler von $Q$ und von $Q'$ (\S.\,4, 5.).
Um $Q$ in die Nebengruppen von $R$ zu zerlegen, wählen wir ein
passendes System von Elementen $1, b_1, b_2, b_3 \ldots$ in $Q$, so
dass $R, Rb_1, Rb_2, Rb_3 \ldots$ alle von einander verschieden
werden und setzen
\[
Q = R + Rb_1 + Rb_2 + Rb_3 + \ldots,
\]
was wir auch, wenn wir
\weq{18}{B = 1, b_1, b_2, b_3 \ldots}
setzen, kurz mit
\weq{19}{Q = RB}
bezeichnen können. Nun ist nach (13) $Q'Q = P$, also auch
\weq{20}{P = Q'RB = Q'B}
oder
\[
P = Q' + Q'b_1 + Q'b_2 + Q'b_3 \ldots
\]

Die Nebengruppen
\weq{21}{Q',\ Q'b_1,\ Q'b_2,\ Q'b_3 \ldots}
sind aber alle von einander verschieden. Denn wäre etwa
\[
Q'b_2 = Q'b_1,
\]
so würde folgen, da $Q'$ die Einheit enthält, dass es ein Element
$e$ in $Q'$ giebt, so dass
\[
b_2 = eb_1
\]
wäre. Dies Element $e$ ist aber gleich $b_2b_1^{-1}$, und also, da
die $b$ in $Q$ enthalten sind, gleichfalls in $Q$ und also auch in
$R$ enthalten. Dann aber wäre auch $Rb_2 = Rb_1$, was gegen die
Voraussetzung ist.

Hiernach können wir die Elemente der zu $R$ complementären Gruppe
in Bezug auf $Q$ und der zu $Q'$ complementären Gruppe in Bezug auf
$P$ bilden. Wir erhalten diese beiden Gruppen
\weq{22}{\begin{aligned}
Q/R &= R,\ Rb_1,\ Rb_2,\ Rb_3 \ldots \\
P/Q' &= Q',\ Q'b_1,\ Q'b_2,\ Q'b_3 \ldots
\end{aligned}}
